% ============================================================================
% Cover Letter — Paper 4 (JCAMD)
% Journal of Computer-Aided Molecular Design (Springer)
% ============================================================================
\documentclass[10pt,a4paper]{letter}

\usepackage[margin=1.7cm]{geometry}
\usepackage{microtype}
\usepackage{siunitx}
\usepackage{hyperref}
\usepackage[nameinlink,noabbrev]{cleveref}
\usepackage{enumitem}

\sisetup{
  detect-all           = true,
  group-minimum-digits = 4,
  range-phrase         = {--},
  range-units          = single,
}

\signature{Myke Vital Sao Temgoua \\ \small Corresponding Author}
\address{%
  University of Yaound\'{e} I \\
  Department of Physics, Faculty of Science \\
  P.O. Box 812, Yaound\'{e}, Cameroon \\
  \href{mailto:myke-vital.sao@facsciences-uy1.cm}{myke-vital.sao@facsciences-uy1.cm}
}

\begin{document}

\begin{letter}{%
  The Editors \\
  \textit{Journal of Computer-Aided Molecular Design} \\
  Springer Nature
}

\opening{Dear Editors,}
\small
\setlength{\parskip}{0.35em}

We are pleased to submit our manuscript, \textbf{``Pareto-guided Monte Carlo tree search for analysing multi-objective alternatives in antimalarial molecular design''}, for consideration as a Research Article in the \textit{Journal of Computer-Aided Molecular Design}.

This work addresses a fundamental challenge in computer-aided molecular design: how to navigate multi-objective trade-offs when potency, synthetic accessibility, and resistance-related criteria cannot be optimally satisfied simultaneously. Our manuscript introduces a scaffold-aware Pareto-guided Monte Carlo tree search (MCTS) for antimalarial fragment assembly that retains chemically distinct candidate profiles rather than collapsing them into a single scalar ranking. We believe this decision-oriented computational framework is well aligned with JCAMD's emphasis on rigorous, transparent computer-based methods for molecular analysis and design.

The study presents two complementary evaluations. First, a \num{20}-seed fixed-budget benchmark quantifies scalar optimisation efficiency against three baselines. The results constitute an honest negative finding: random search achieved the highest mean reward (\num{0.6724} $\pm$ \num{0.0056}), followed by MCTS (\num{0.6649} $\pm$ \num{0.0068}), genetic algorithm (\num{0.6453} $\pm$ \num{0.0124}), and greedy search (\num{0.4278}). MCTS incurs a measurable scalar-performance cost ($t_{19} = -4.97$, $p = 0.000085$) relative to broad random exploration in this relatively flat reward landscape. Second, a separate Pareto analysis identifies four non-dominated candidate profiles (hypervolume \num{1.2366}) that expose potency--accessibility trade-offs invisible to any fixed scalar weighting. These profiles differ in resistance-informed chemotype similarity and polypharmacology-informed docking scores—computational design proxies that inform decision-making but do not substitute for experimental validation.

The manuscript's central methodological contribution lies in separating optimisation efficiency from candidate-set geometry. The scalar benchmark quantifies the computational cost of maintaining a Pareto archive; the candidate-level analysis demonstrates the decision support value that archive provides. This dual evaluation reflects the journal's standards for transparent computational validation: we report what the algorithm optimises well (Pareto front diversity, scaffold-aware chemistry) and what it does not (scalar reward maximisation under this specific objective function and search budget). The resistance-informed term employs Morgan/Tanimoto similarity to reference chemotypes from Project 2, and the polypharmacology-informed term aggregates multi-target docking scores; these are explicitly presented as computational proxies, not direct biological measurements.

Key features aligned with JCAMD's scope include: (1) a chemistry-informed PUCT prior that focuses expansion on synthetically valid fragment combinations, reducing wasted oracle calls on chemically unproductive branches; (2) transparent reporting of an honest-negative scalar benchmark alongside the Pareto archive analysis, demonstrating that methodological value need not require algorithmic superiority claims; (3) rigorous statistical analysis (paired $t$-tests, $2^5$ factorial ablation with main effects) quantifying the contributions of scaffold-aware policies ($\Delta = +0.148$) and Pareto maintenance ($\Delta = +0.108$); and (4) full computational reproducibility, with public code, molecule sets, and analysis outputs available under the MIT licence at \url{https://github.com/NanaEngo/Malaria_codesV2}.

This work is particularly relevant to JCAMD readers working on de novo molecular generation, multi-objective optimisation in drug discovery, and computational methods for neglected tropical diseases. The fragment-assembly framework is general and could be adapted to other therapeutic targets where decision-makers need to evaluate candidates under different weighting schemes—for example, early-stage discovery favouring novelty versus late-stage optimisation prioritising synthetic feasibility.

This work is original, is not under consideration elsewhere, all authors have approved the submission, and the authors declare no competing financial interests. We are submitting under the traditional subscription publishing model.

\closing{We thank you for your consideration and look forward to the reviewers' feedback.}

\ps{%
\textbf{Suggested reviewers:}\\
\begin{itemize}[leftmargin=*, itemsep=0pt, topsep=0pt]
  \item Dr.\ Jan H.\ Jensen (University of Copenhagen) — graph-based genetic algorithms and Monte Carlo methods for chemical space exploration
  \item Prof.\ Gisbert Schneider (ETH Zurich) — generative molecular design and de novo drug discovery methodologies
  \item Dr.\ Ola Engkvist (AstraZeneca / Chalmers University) — multi-objective optimisation and AI-driven molecular design
  \item Prof.\ Artem Cherkasov (University of British Columbia) — computer-aided drug design for neglected tropical diseases including malaria
\end{itemize}
}

\end{letter}
\end{document}
